%------------------------------------------------------------------------------
% Template dos artigos da Revista LUPS
%------------------------------------------------------------------------------

\documentclass[12pt,a4paper,compsoc]{IEEEtran}
\usepackage[utf8]{inputenc}
\usepackage[brazilian]{babel}
\usepackage[dvips]{graphicx}
\usepackage{epsfig}
\usepackage{amsfonts, color}
\usepackage{subfigure}
\usepackage{eurosym}
\usepackage[T1]{fontenc} 
\usepackage{ae}
\usepackage{cite}

%------------------------------------------------------------------------------

\begin{document}

%------------------------------------------------------------------------------
%%%%%%%%%%%%%%%%%%%%%%%%%%%%%%%%%%%%%%%%%%%%%%%%%%%%%%%%%%%%%%%%%%%%%%%%%%%%%%%
% Estas informaes devem ser alteradas pelo editor da revista
\pubid{LABORATORY OF UBIQUITOUS AND PARALLEL SYSTEMS~\copyright~LUPS} 
\renewcommand{\leftmark}{REVISTA~LUPS,~VOL.~2, NO.~1,~DEZEMBRO~2013}
%%%%%%%%%%%%%%%%%%%%%%%%%%%%%%%%%%%%%%%%%%%%%%%%%%%%%%%%%%%%%%%%%%%%%%%%%%%%%%%
%------------------------------------------------------------------------------

% Ttulo e autores do artigo

\title{Explorando Regras para Consciência de Situação em Ubicomp}

\author{Caue Duarte, UFPel; Alexandre Gomes da Costa, UFPel

\IEEEcompsocitemizethanks{\IEEEcompsocthanksitem \textbf{Caue Duarte}: Programa de Pós-Graduao em Computação,  Universidade Federal de Pelotas - UFPel, Centro de Desenvolvimento Tecnológigo - CDTec.
% Para uma quebra de linha, deve-se utilizar o \protect antes do \\, caso contrrio vai dar erro.
\protect\\ E-mail: e-mailautor1@inf.ufpel.edu.br}

% No caso de mais autores, basta repetir:
\IEEEcompsocitemizethanks{\IEEEcompsocthanksitem \textbf{Alexandre Gomes da Costa}: Programa de Pós-Graduao em Computação,  Universidade Federal de Pelotas - UFPel, Centro de Desenvolvimento Tecnológigo - CDTec.
\protect\\ E-mail: alexandre.costa@inf.ufpel.edu.br}
}

%--------------------------------------------------------------------------------------------------
% ABSTRACT
%--------------------------------------------------------------------------------------------------
\IEEEcompsoctitleabstractindextext{%
\begin{abstract}
  Técnicas de Consciência de Situação são amplamente utilizadas por diversas áreas dentro das
  Ciências da Computação, normalmente ligadas a consequência dos estudos em consciência de
  contexto, com o diferencial possuir raciocínio agregado. Nesse contexto, buscamos no presente
  trabalho definir principais conceitos referentes a computação ubíqua e consciência de situação
  além de explorar o emprego de sistemas baseados em regras para prover esse raciocínio, explorando
  o emprego de regras em trabalhos de grande porte em computação ubíqua.
\end{abstract}

\begin{IEEEkeywords}
Consciência de Situação, Sistemas baseados em regras, Ubicomp, Computação Pervasiva.
\end{IEEEkeywords}}

\maketitle

%--------------------------------------------------------------------------------------------------
% INTRODUÇÃO
%--------------------------------------------------------------------------------------------------
\section{Introdução}

  \IEEEPARstart{C}{}om o crescente número de dispositivos móveis e o aumento significativo de seu
  poder de processamento e largura de banda, o paradigma da computação ubíqua proposta por Weiser,
  que previa um aumento nas funcionalidades e na disponibilidade de serviços de computação para os
  usuários finais com um custo de gerenciamento ficando cada vez menor, torna-se cada vez mais
  viável.
  
  Técnicas de Consciência de Situação são amplamente utilizadas por diversas áreas dentro da
  Ciência da Computação, incluindo a Computação Ubíqua. Normalmente ligadas a consequência dos
  estudos em consciência de contexto, com o diferencial possuir raciocínio agregado.
  
  Este trabalho explora o uso de Sistemas Baseados em Regras para prover Consciência de Situação.
  Regras podem ser uma generalização geralmente válida, um princípio regulador, um método para a
  realização de uma determinada operação matemática e obter um certo resultado.
  
  Foram analisados os seguintes projetos Collaborative Context-Aware Service Platform for Pervasive
  Computing (CoCA), WComp, a Middleware for Ubiquitous Computing (WComp), A Rule-Based Platform for
  Situation Management (SCENE) e por fim Framework and Rule- based Language for Facilitating 
  Context-aware Computing using Information Appliances - (Context-Aware rule Description Language -
  CADEL).
  
  O objetivo deste trabalho é apresentar conceitos relacionados a Computação Ubíqua, Consciência de
  Situação, Sistemas Baseados em Regras além de explorar o uso de Regras para prover o raciocínio
  utilizado em Consciência de Situação.
  
  O artigo é descrito como segue: na seção 2 são apresentados conceitos de Computação Ubíqua. A
  seção 3 apresenta um breve resumo de Conciência de Situação. A seção 4 introduz conceitos de
  Sistemas Baseados em Regras e mostra sua utilização em projetos relacionados e finalmente na
  seção 5 são apresentadas as conclusões do trabalho.


%--------------------------------------------------------------------------------------------------
% COMPUTAÇÃO UBÍQUA
%--------------------------------------------------------------------------------------------------
\section{Computação Ubíqua}

  O nome Ubíquo vem do Latim \textit{ubiquu}, que significa estar em todos os locais. Mark Weiser,
  então cientista do Centro de Pesquisa da Xerox, definiu em seu artigo intitulado  \textit{The
  Computer for the 21st Century} o termo Computação Ubíqua como sendo a presença direta e constante
  de serviços de computação na vida dos usuários finais, dividindo a computação entre diversos
  dispositivos conectados entre si. O desafio dos cientistas é tornar real o pensamento de Weiser,
  e consequentemente a ubiquidade na computação, fazendo com que as interfaces se tornem
  transparentes, armazenem informações e aprendam apartir dessas informações. Compartilhar dados,
  informações de forma mais fácil e transparente, esta tecnologia fará parte do processo de
  evolução do ser humano, onde ca vez mais a computação centralizada deixar de ser o único agente
  que o usuário interage e multiplicar-se para usos especializados. A Computação Ubíqua,  não se
  resume apenas pela computação, seu objetivo é integrar totalmente a relação
  computação/periféricos com o usuário tornando-se  invisível, no sentido de a utilizar sem 
  perceber.
  
  A computação, embarcada nos mais diversos computadores, já faz parte da vida das pessoas em maior
  ou menor grau e essa interação juntamente com o número da computadores tente a aumentar,
  tornando-se onipresente na vida do usuário. Para tornar isso possível, busca-se a utilização de
  interfaces naturais, tornando a comunicação com esses dispositivos mais transparente, interagindo
  com o usuário por meio de gestos, fala e visão, por exmeplo.
  
  Outra forma de interação é através da computação sensível ao contexto, que através do 
  processamento de dados oriundos dos mais diversos sensores, percebe e interage com a realidade do
  ambiente como por exemplo a movimentação do usuário através de um ambiente monitorado.
  O computador, presente nos mais diversos dispositivos (smartphones, tablets, etc..), possui 
  interfaces de comunicação, e como todas interfaces já consolidadas, tende a se tornar 
  imperceptivel, uma extensão dos sentidos humanos. Como exemplo de interfaces podemos citar um
  óculos ou uma bengala, que pelo uso contínuo, o usuário abstrai o utensilio em si, utilizando-o
  de forma transparente, focando apenas na realidade trazida pelo mesmo. Sob esta premissa, o uso
  das interfaces passa a tornar-se inconsciente~\cite{weiser1993}. Nesse contexto as aplicações 
  precisam ``entender'' e se ajustar ao ambiente, através do processamento de informações sobre o
  contexto em que estão inseridas ~\cite{maciel2004}.
  
  Essa nova classe de sistemas computacionais, conscientes ao contexto, abre perspectivas para o
  desenvolvimento de aplicações muito mais ricas, elaboradas e complexas, que exploram a natureza
  dinâmica e a mobilidade do usuário. Um desafio central na programação deste tipo de aplicação é
  possibilitar que as mesmas se adaptem continuamente ao ambiente e permaneçam funcionando mesmo
  quando o indivíduo se movimentar ou trocar de dispositivo ~\cite{GrimmB03, Cac08}.
  
  A Computação Ubíqua é entendida como o terceiro grande paradigma computacional, precedido pelos
  \textit{mainframes} e pela computação pessoal ~\cite{weiser1997}. Com foco em demandas 
  institucionais, hoje atende demandas específicas, os \textit{mainframes}, computadores de grande
  porte, atendem requisições de inúmeros usuários. Já com a computação pessoal, observamos uma
  convergência das máquinas para uso pessoal.
  
  Visualizamos a computação ubíqua como uma rede invisível que permeia o usuário com dispositivos
  inteligentes e que se intercomunicam, aprendem uns com os outros, com o ambiente e tomam
  decisões. Os recusos computacionais devem estar disponíveis ao usuário, promovendo adaptações
  dinâmicas tanto dos serviços quanto aplicações satisfazendo as necessidades do usuário, de 
  maneira mais transparente possível. Considerando esta perspectiva, a Computação Ubíqua constitui
  um ambiente altamente distribuído, heterogêneo, dinâmico, móvel, mutável e com forte interação
  entre homem e máquina ~\cite{Augustin03}.
  
  Um caminho breve para a Computação Ubíqua é a imersão total do usuário nessa rede invisível, não
  só de computadores, \textit{smartphones}, \textit{tablets}, mas de dispositivos comuns, até então
  não associados com computação, como eletrodomésticos ao exemplo das já consolidadas 
  \textit{smartv's}, além de outros tantos como mesas, talheres e até ambientes inteiros. Através
  de estratégias de \textit{cloud computing} (computação em nuvem) interligando esses dispositivos
  por rede, internet, \textit{bluetooth} podem trocar dados, processar informações e executar
  qualquer ação que possa ser interessante do ponto de vista do usuário ou sistema.
  

%------------------------------------------------------------------------------

\section{Projetos em Conscincia de Situao Explorando Regras}

Um dos primeiros e, portanto, mais amplamente desenvolvidos modelos de tomada de decisão são sistemas de baseados regras. Seu  desenvolvimento começou na década de 60, mas tornou-se comum nos anos 70 e 80.

Regras podem ser  uma generalização geralmente válida , um princípio regulador , um método para a realização de uma determinada  operação matemática e obter um certo resultado. Exemplos do uso de regras  :

Regras em sistemas especialistas

IF (condição 1) AND (condição 2) AND (condição 3) AND (condição 4) THEN (solução 1)


Regras na Web Semântica utilizando SWRL:

condição1(?x1,?x2) $\wedge$ condição2(?x2,?x3) $\Rightarrow$ solução1(?x1,?x3)

Regras na Computação Ubíqua  em  ECA-DL


Sobre  Regra1(p)

Quando Condição (p)

Faça Ação  (p, “mensagem”)


Regras em bancos de dados ( regras  de consistência)

CREATE TABLE acidtest (A INTEGER, B INTEGER CHECK (A + B = 100));


Interfaces de configuraçaõ de regras de alto nível

(retirar imagem da tela de configuração do cobalto , ou gmail,) 


Sua representação assume a forma de pares antecedente-consequente ou instruções if-then, Figura 2.18;



Além de suporte para a lógica fuzzy, o método é diferente em termos de,

\begin{enumerate}
\item Apenas uma regra fica para fornecer a ação;
\item Arbitragem necessária para determinar qual regra vencerá
\end{enumerate}




As vantagens do uso de regras são:

A existência de base teórica.
Possibilidade de reutilização de regras
Simplicidade de projeto. 
Permite a separação das regras do resto do sistema.
Utilização de plataformas de regras .
Permite adaptação de sistemas on-the-fly
Entrega incremental de sistemas
Explanação do processo de inferência
Permite gerar conclusões com grau de incerteza (a partir de premissas imprecisas)


Algumas desvantagens são :

Desempenho
Pouco relacionamento entre regras (regras se relacionam através do processo de inferência apenas)
Podem gerar erros
Dificuldade de lidar com ambiguidades
Regras conflitantes
Base de conhecimento  pode ficar inconsistente e/ou desatualizada

\textbf{Arquitetura de Sistemas Baseados em Regras}

Cinco componentes básicos compreendem um sistema baseado em regras genéricas,

\begin{enumerate}

\item Banco de dados de Matching: 

O banco de dados corresponde ao conjunto de recursos (premissas) que a base de regra usa para estabelecer estados. Os recursos são avaliados periodicamente para quaisquer situação contra o conteúdo das regras expressas no sistema baseado em regras .

\item Regras de Condição (Condition-Action Rules) Uma vez estabelecido o banco de dados de fatos a respeito do status dos dados, a regra de condição será expressa conforme exemplo:

\begin{itemize}
\item IF (condição 1) AND (condição 2) THEN (ação)
\end{itemize}

\item Reescrevendo Regras:

Podem existir situações em que uma base de dados é alterada diretamente, sem mais raciocínios, tal comportamento é autônomo e geralmente reflete comportamentos NPC mais reativos.
Tais casos s ão expressos como regras de reescrita 

\begin{itemize}
\item IF (condição 1 satisfeita) AND (condição 2 satisfeita) THEN 
\item remove(ação 1)
\item add(ação 2 )
\end{itemize}

\item Encadeamento 

O encadeamento  começa com o conhecimento do objetivo  representado, e em seguida tenta encontrar um caminho mínimo através da base de regras para os estados registrados no banco de dados .

 
\item Arbitragem: 
Geralmente limitar o número de regras para um único caso. Uma das várias heurísticas são geralmente empregadas para esse fim.

\end{enumerate}


\subsection{CoCA - Hybrid Approach to Collaborative Context-Aware Service Platform for Pervasive Computing}

CoCA É um middleware baseado em vizinhança, que visa a aquisição e utilização de informações de contexto para fornecer serviços, com um modelo de gerenciamento de contexto híbrido de grande escalabilidade, formalidade e reutilização.
O Projeto CoCA [citar] apresenta uma plataforma colaborativa serviço conscientes de contexto baseado no modelo de gestão de contexto híbrido. 
O comportamento das aplicações difundidas não dependem apenas em suas interações estaduais e de usuários internos, mas também na contextos detectado durante a sua execução. Realizando raciocínio e decisões com base em dados de contexto, semântica contexto e as regras e políticas relacionadas. Tais dados são organizados em um modelo de gestão contexto híbrido, chamado HCOM [citar] . Mais tarde, atualizando a ontologia com base em modelo de gestão de contexto genérico surge o, GCOM [citar].

Apoio à decisão é fornecido de forma proativa ou reativa. É fornecido um suporte à decisão e funcionalidade de ação para minimizar a intervenção do usuário.

IMAGEM Fig.5
%\begin{figure*}[t!]
%\centerline{\includegraphics[scale=.2]{coca.png}}
%\caption{Component view of the Enhanced CoCA platform}
%\label{coca}
%\end{figure*}

%Figura~\ref{coca}

\subsection{WComp, a Middleware for Ubiquitous Computing}

WComp, é um middleware para a computação ubíqua. Esse modelo de Middleware é baseado em três partes distintas: uma infra-estrutura de software,  um serviço de composição da arquitetura e um mecanismo de adaptação.

Para gerenciar o dinamismo e a heterogeneidade das entidades na infra-estrutura de software é utilizado WSOAD ( Web Service Oriented Architecture for Device). As aplicações ubíquas são baseadas em um conjunto de serviços Web para dispositivos. Os serviços não são editáveis, logo para adicionar novas funcionalidades, uma aplicação tem que ser uma composição de serviços para dispositivos.

O WComp utiliza o conceito de Aspectos da Assembléia (AA), baseado em programação orientada a aspectos. Eles definem algumas reconfigurações estruturais de um aplicativo que são acionados em resposta para eventos da infra-estrutura de software. Eventos informam sobre o aparecimento ou desaparecimento de dispositivos da infra-estrutura de software. 

Estas regras são compostas e criadas com uma lógica bem definida, no caso de conflito. Eles são aplicados em conjuntos de componentes que não são necessariamente conhecidas a priori. Várias linguagens e composição de regras podem ser definidas de acordo com o tipo de aplicações nas quais eles serão aplicados.


\subsection{SCENE - A Rule-Based Platform for Situation Management}

Nesse artigo foi proposto uma abordagem para a especificação e realização de detecção de situação para uma aplicação de consciência de situação. Foi implementado uma plataforma baseada em regras para a gestão de situação (SCENE) que aproveita JBoss Drools adicionando funcionalidade para suportar nativamente consciência de situação baseado em regras.
As regras são definidas no Drools por meio de uma linguagem específica de domínio chamado Drools Rule Language (DRL). A declaração regra DRL compreende uma condição e um bloco de expressão conseqüêntes, respectivamente conhecido como  Left Hand Side (LHS) e Right Hand Side (RHS). A regra especifica que quando o determinado conjunto de condições definidas na LHS ocorre, a lista de ações na RHS deve ser executado. A LHS é composta de elementos condicionais que podem ser combinadas através de operadores lógicos e operadores SET, como  contains e member of. A RHS permite a declaração do código que vai ser executado quando as condições definidas na LHS forem satisfeitas.

Essas características de situações nos levam aos seguintes requisitos básicos para a nossa abordagem de situação:
Tipos de situação deve ser definida em tempo de design, e situações instanciar esses tipos devem ser detectados em tempo de execução; 
Tipos de situação deve ser definido com referência a tipos de entidades, bem como restrições sobre propriedades e relações das entidades;
Propriedades temporais das situações devem ser considerados (tais como o tempo inicial, e, para uma situação passado, o tempo final e duração).
A fim de atender a estes requisitos, tipos de situação são especificadas no SCENE por meio de aspectos estruturais e comportamentais, que são representados por  Situation Classes e Situation Rules, respectivamente. Cada Situation Class definido pelo usuário é a especialização da classe pré-definida SituationType, que é uma classe abstrata que trata das propriedades temporais da situação e características de composição. A FIGURA X ilustra uma Situation Class, Fever, que extends uma classe SituationType.
Figura(Situation Classe - fever.png)

Nesse trabalho eles encaram situações como sendo entidades compostas cujos constituintes são outras entidades, suas propriedades e as relações em que estão envolvidos. Como exemplo de situação podemos ter: "John is working", "John has fever", "John has had an intermittent fever for the past 6 months", "John and Paul are outdoors, at a distance of less than 10m from each other", "Bank account number 87346-0 is overdrawn while a suspicious transaction is ongoing", etc.


Especificar as situação exige um padrão de regra única seguido de um dialeto de restrição de regras padrão Drools. Situações pode ser composta de restrições sobre domínio de entidades, e além disso, pode ser composto por situações existentes nelas mesma. Nesse trabalho foram abordados aspectos temporais das aplicações, e foi incluído operadores para relacionar situações baseadas em seus aspectos temporais. A detecção é baseada em regras, e é implantado um mecanismo de regras maduro e eficiente e uma complexa tecnologia de processamento de eventos disponíveis para uso. A plataforma gerencia situações através da implementação de situações de controle do ciclo de vida, como a ativação da situação, a manutenção do estado e desativação. Não foi avaliado o desempenho da detecção da situação, mas esta em curso e por experiência em trabalhos anteriores é esperado que o desempenho da detecção de situação seja adequada para a maioria das aplicações. 


\subsection{CADEL - Framework and Rule-based Language for Facilitating Context-aware Computing using Information Appliances}

Recentemente, computação consciente do contexto, utilizando aparelhos de informação é um dos tópicos de pesquisa mais importantes. Em sistemas de computação consciente do contexto, os dispositivos são controlados automaticamente com base no contexto atual obtido a partir de vários sensores, como posição do usuário, temperatura ambiente e assim por diante. A fim de fazer com que os sistemas sensíveis ao contexto funcionem corretamente, é preciso identificar o contexto atual do ambiente (incluindo usuários) e para recuperar as regras que podem ser executadas no contexto. Como técnicas para descobrir os dispositivos específicos em um ambiente ubíquo. Mesmo com essas técnicas, a fim de fazer vários dispositivos trabalhar cooperativamente com base no contexto, um cenário de como usuário(s) deseja controlar esses dispositivos, deve ser especificado com antecedência. Em um cenário, as regras que consistem em condições de sensores e ações de dispositivos são descritos. Para especificar cada regra corretamente, é necessário escolher com cuidado um conjunto de sensores e dispositivos que o usuário deseja controlar. Assim, os usuários têm que estar familiarizados com as suas funcionalidades. No entanto, seria muito difícil para os usuários em uma casa comum descrever um cenário viável para fazer o sistema funcionar em seus caminhos esperados. Além disso, em um ambiente familiar, vários usuários podem querer controlar o mesmo dispositivo em simultâneo de diferentes maneiras.

Neste trabalho, foi proposto um quadro para permitir que todos possam facilmente descrever cenários para sistemas de computação consciente do contexto, incluindo vários aparelhos de informação e sensores. Nossa estrutura facilita (1) personalização de dispositivos, (2) especificação intuitivo de regras, e (3) verificação de consistência e de conflito detecção em várias regras.

Para os fins acima, em primeiro lugar, foi definido uma linguagem chamada CADEL (Context-Aware Linguagem de descrição de regras) para especificar regras. Desde CADEL tem sintaxe e semântica para linguagens naturais semelhantes, usuários domésticos comuns pode especificar regras de forma intuitiva. Em CADEL, cada usuário pode definir novas palavras (por exemplo, hot-and-entupido) para indicar contextos específicos detectados a partir de vários sensores.

Em segundo lugar, a estrutura oferece uma função de orientação aos usuários durante a descrição da regra, com o qual os usuários podem recuperar os sensores e dispositivos próximos através de GUI e obter as informações, tais como as ações permitidas de um dispositivo e o valor de um sensor. Consequentemente, os usuários podem facilmente especificar uma regra de utilização da informação obtida.

Em terceiro lugar, o quadro fornece um mecanismo para detectar automaticamente um conflito entre várias regras, o que acontece quando as condições de várias regras mantidas ao mesmo tempo e executam ações diferentes para o mesmo dispositivo. A fim de evitar inconsistências devido a um conflito, foi especificado a prioridade entre essas regras conflitantes. Quando uma nova regra é registrada, o quadro verifica se há conflitos de regras com as regras vigentes. Se ele entra em conflito, o quadro solicita aos usuários para especificar a prioridade entre as regras. Os usuários podem anexar um contexto específico para a prioridade de modo que a prioridade só funciona no contexto.

Tem sido implementado um protótipo de sistema do quadro proposto em um PC usando UPnP, que avaliou o desempenho do sistema em termos de tempo de resposta de recuperação do sensor/dispositivo e detecção de conflitos ao longo de várias regras. Através de experimentos, foi confirmamos que a implementação do protótipo alcança um desempenho praticamente suficiente para executar essas operações.


%------------------------------------------------------------------------------

\section{Consideraes Finais}

Apresentar as considerações finais do capítulo.

%------------------------------------------------------------------------------

\bibliographystyle{unsrt}
\bibliography{BIBFile}

%------------------------------------------------------------------------------
%Biografia dos autores

\begin{IEEEbiography}[{\includegraphics[width=1in, height=1.25in, clip, keepaspectratio]{imagens/fotodoautor}}]{Caue Duarte}
Possui graduação em Análise e Desenvolvimento de Sistemas pela Universidade Norte do Paraná. Técnico em Eletrônica pelo Instituto Federal Sul-Rio-grandense. Possui Pós-Graduação com Enfase em Educação à Distância , atualmente é mestrando no Programa de Pós-Graduação em Computação da UFPel (Mestrado em Ciência da Computação). Atua como técnico em Sistemas de Informação da Universidade Federal de Pelotas, onde desenvolve sistemas, interfaces e realiza modelagem de sistemas/banco de dados e estudo de casos e é professor do Instituto Educacional Dimensão, no curso de Análise de Sistemas. 
%Quanto a foto, no deve ser de corpo inteiro (deve ser similar ao avatar de exemplo) e convertida em escala de cinza.
\end{IEEEbiography}

\begin{IEEEbiography}[{\includegraphics[width=1in, height=1.25in, clip, keepaspectratio]{imagens/fotodoautor.png}}]{Alexandre Gomes da Costa}
Aluno Especial do Programa de Ps-Graduao em Computao da Universidade Federal de Pelotas. Possui Bacharelado em Cincia da Computao (2008) pela Universidade Federal de Pelotas.
%Quanto a foto, no deve ser de corpo inteiro (deve ser similar ao avatar de exemplo) e convertida em escala de cinza.
\end{IEEEbiography}

\end{document}


